\documentclass[12pt,a4paper,oneside]{report}
\usepackage[utf8]{inputenc}
\usepackage[english]{babel}
\usepackage{geometry}
\geometry{a4paper, margin=1in}
\usepackage{setspace}
\onehalfspacing 
\usepackage{booktabs}
\usepackage{titlesec}

\titleformat{\chapter}[display]
  {\normalfont\bfseries\huge}{\chaptername\ \thechapter}{20pt}{\Huge}
\titleformat{\section}
  {\normalfont\selectfont\bfseries\large}{\thesection}{1em}{}

\begin{document}

\setcounter{chapter}{0} 

\clearpage
\thispagestyle{empty} 
\vspace*{\fill}
\begin{center}
    {\Huge\bfseries CHAPTER 1} \\[0.5cm]
    {\Huge\bfseries INTRODUCTION}
\end{center}
\vspace*{\fill}
\clearpage

\chapter{Introduction}

\section{Background of the Study}
In the era of Industry 4.0, Information Technology Service Management (ITSM) has become the backbone of institutional efficiency and digital transformation. The ITIL 4 framework provides a globally recognized approach for creating value through service management. However, as organizations migrate to complex cloud environments, traditional service operations face significant challenges in agility and decision-making. Emerging technologies, particularly Generative AI (GenAI) and Robotic Process Automation (RPA), offer promising solutions to transform manual workflows into ``Intelligent Governance'' models.

\section{Problem Statement}
Despite the theoretical benefits of ITIL 4, many academic and governmental institutions in Yemen, such as Sana'a University, continue to rely on manual processes for monitoring and reporting Key Performance Indicators (KPIs). This reliance leads to a significant ``meaning gap'' between raw technical data and strategic institutional objectives. Current literature shows a lack of empirical frameworks that integrate GenAI and RPA for automated KPI assessment. Furthermore, there is a limited application of multi-criteria decision-making methods, such as the SECA (Simultaneous Evaluation of Criteria and Alternatives) method, to objectively weight and select KPIs in resource-constrained environments. Without an automated and intelligent approach, ITSM remains reactive rather than proactive, hindering the efficiency of digital services.

\section{Purpose of the Study}
The purpose of this quantitative study is to investigate the effectiveness of an intelligent framework that integrates GenAI and RPA into ITIL 4 processes to improve the selection, monitoring, and assessment of KPIs at the Faculty of Computing and IT (FCIT), Sana'a University.

\section{Research Questions}
The study addresses the following research questions to bridge the identified gaps:
\begin{itemize}
    \item \textbf{Main Question:} How can the integration of AI and RPA enhance the automation and accuracy of ITIL 4 Key Performance Indicators?
    \item \textbf{Sub-Questions:}
    \begin{itemize}
        \item What are the primary ITIL 4 metrics in service operations that can be optimized using GenAI and RPA?
        \item To what extent does the application of the SECA method improve the objective weighting and selection of KPIs compared to traditional methods?
        \item What are the technical and organizational challenges of implementing automated ITSM governance in Yemeni higher education institutions?
    \end{itemize}
\end{itemize}

\section{Research Objectives}
To achieve the purpose of this study, the following objectives have been formulated:
\begin{enumerate}
    \item To analyze the current challenges and limitations of manual KPI monitoring within the ITSM practices at Sana'a University.
    \item To design a prototype framework for automated ITIL service operations leveraging RPA and AI-driven insights.
    \item To evaluate the performance of the proposed framework using the SECA method for multi-criteria decision-making.
    \item To provide recommendations for adopting ``Intelligent Governance'' in the Yemeni academic context.
\end{enumerate}

\section{Significance of the Study}
\begin{itemize}
    \item \textbf{Theoretical Significance:} This study contributes to the body of knowledge by bridging the gap between ITIL 4 governance and intelligent automation (AI/RPA) in developing-country contexts.
    \item \textbf{Practical Significance:} It provides a practical roadmap for IT managers at Sana'a University to reduce incident resolution time, minimize human error, and align IT services with academic goals.
    \item \textbf{Methodological Significance:} It demonstrates the utility of the SECA method as a robust tool for objective performance measurement in computing environments.
\end{itemize}

\section{Scope and Delimitations}
The study is delimited to the ITIL 4 Service Operations processes (specifically Incident and Problem Management) at the Faculty of Computing and IT (FCIT), Sana'a University. The technical scope focuses on the integration of GenAI and RPA tools. The timeframe for this research is the academic year 2025-2026.

\section{Limitations}
Potential limitations include the restricted access to high-performance computing resources for AI training and the possible scarcity of comprehensive historical ITSM logs in some departments, which may affect the initial data analysis phase.

\section{Definition of Key Terms}
To ensure clarity and consistency throughout this study, the following operational definitions are provided:
\begin{itemize}
    \item \textbf{ITIL 4:} The latest evolution of the Information Technology Infrastructure Library framework, focusing on the Service Value System (SVS) and value co-creation.
    \item \textbf{Key Performance Indicators (KPIs):} Quantifiable metrics used to evaluate the success of IT service operations, such as incident resolution time and system availability.
    \item \textbf{Robotic Process Automation (RPA):} Software technology used to automate repetitive and rule-based tasks within ITIL processes, such as data extraction and reporting.
    \item \textbf{Generative AI (GenAI):} A branch of artificial intelligence used in this study to classify IT incidents and provide intelligent insights for automated decision-making.
    \item \textbf{SECA Method:} (Simultaneous Evaluation of Criteria and Alternatives) A multi-objective optimization and decision-making mathematical model used to objectively weight and rank the selected KPIs.
\end{itemize}

\section{Organization of the Thesis}
This thesis is organized into five interrelated chapters to ensure a logical flow of the research:
\begin{itemize}
    \item \textbf{Chapter 1 (Introduction):} Provides the foundation of the study, including the background, problem statement, purpose, research questions, objectives, significance, scope, and definition of key terms.
    \item \textbf{Chapter 2 (Literature Review):} Reviews the existing body of knowledge regarding ITIL 4 framework, the role of AI and RPA in ITSM, and contemporary methods for KPI measurement.
    \item \textbf{Chapter 3 (Research Methodology):} Details the quantitative research design, the data collection procedures, and the application of the SECA method for performance assessment.
    \item \textbf{Chapter 4 (Implementation and Results):} Presents the development of the automated framework and discusses the findings derived from the data analysis.
    \item \textbf{Chapter 5 (Conclusion and Recommendations):} Summarizes the study’s contributions, draws final conclusions, and suggests directions for future research in the field of intelligent IT governance.
\end{itemize}

\section{Research Timeline}
The following timeline outlines the expected duration for each phase of the research, ensuring the study is completed within the academic year 2025-2026:

\begin{table}[h!]
\centering
\caption{Research Timeline and Key Phases}
\begin{tabular}{lll}
\toprule
\textbf{Phase} & \textbf{Research Activity} & \textbf{Duration} \\
\midrule
Phase 1 & Literature Review \& Finalizing Proposal & Month 1 - 2 \\
Phase 2 & Data Collection (ITSM Logs \& Expert Surveys) & Month 3 \\
Phase 3 & Framework Design (RPA Workflows \& AI Integration) & Month 4 - 5 \\
Phase 4 & Data Analysis \& SECA Model Implementation & Month 6 - 7 \\
Phase 5 & Results Evaluation \& Thesis Writing & Month 8 \\
Phase 6 & Final Review \& Thesis Defense & Month 9 \\
\bottomrule
\end{tabular}
\end{table}

\end{document}
